\chapter{Applications: Leontief, OLS, Markov, Portfolios}

Almost every linear-algebraic tool from the previous chapters shows up in a
standard economics course. This chapter collects the headline applications:
the Leontief input-output model (an inverse problem), the ordinary least
squares (OLS) estimator (matrix calculus meets the inverse), Hessian
classification of critical points (Sylvester's criterion in action), Markov
transition matrices (eigenvalues with $|\lambda| \le 1$), and portfolio
variance (a quadratic form). Treat the problems below as integration
exercises drawing on everything so far.

\section{Core Rules}

\begin{rulebox}{Leontief input-output}
An economy with $n$ sectors has a technology matrix $A \in \R^{n \times n}$
where $a_{ij}$ is the amount of good $i$ needed to produce one unit of good
$j$. With final demand $d \in \R^{n}$ the gross output $x$ must satisfy
\[
x \;=\; A x + d \quad \Longleftrightarrow \quad (I - A) x \;=\; d.
\]
If $I - A$ is invertible (which it is when $A$ has spectral radius $< 1$, a
productivity condition), the solution is $x = (I - A)^{-1} d$. The matrix
$(I - A)^{-1}$ is the \emph{Leontief inverse}; its $(i, j)$-entry is total
(direct $+$ indirect) demand on sector $i$ per unit of final demand for
sector $j$.
\end{rulebox}

\begin{rulebox}{OLS estimator}
Given a design matrix $X \in \R^{n \times k}$ of full column rank and
observations $y \in \R^{n}$, the ordinary least squares estimator minimising
$\|y - X \beta\|^{2}$ is
\[
\hat{\beta} \;=\; (X^{T} X)^{-1} X^{T} y.
\]
The $k \times k$ matrix $X^{T} X$ is symmetric and pd when $X$ has full
column rank, hence invertible.
\end{rulebox}

\begin{rulebox}{Hessian classification}
For $f \in C^{2}$ at an interior critical point $x^{*}$, the Hessian
$H = H_{f}(x^{*})$ classifies via the rules of Chapter 7.
\end{rulebox}

\begin{rulebox}{Markov transition matrices}
A column-stochastic matrix $P$ has nonnegative entries with each column
summing to $1$. The state distribution $\pi_{t}$ evolves as
$\pi_{t + 1} = P \pi_{t}$. The eigenvalue $\lambda = 1$ always exists; the
corresponding eigenvector (suitably normalised so its components sum to $1$)
is a stationary distribution.
\end{rulebox}

\begin{rulebox}{Portfolio variance}
For asset weights $w \in \R^{n}$ summing to $1$ and a (symmetric, psd)
return covariance matrix $\Sigma$, the portfolio's return variance is
\[
\mathrm{Var}(w^{T} r) \;=\; w^{T} \Sigma w.
\]
Minimising $w^{T} \Sigma w$ subject to $w^{T} \mathbf{1} = 1$ gives the
global minimum variance portfolio.
\end{rulebox}

\section{Worked Example}

\begin{worked}{A two-sector Leontief economy with $A=\begin{pmatrix}0.2&0.3\\0.4&0.1\end{pmatrix}$ and final demand $\mathbf d=(10,20)^{\!\top}$: find gross output, then check productivity and read the multipliers}
\wstep{Step 1. Know what the entries mean before computing anything.}
$a_{ij}$ is the amount of good~$i$ needed to produce \emph{one unit} of
good~$j$. So $a_{12}=0.3$ says sector~$2$ consumes $0.3$ units of good~$1$
per unit of its own output. Columns are recipes; rows are deliveries.
\wtrap{Reading $A$ transposed is the most expensive error in the chapter,
because the arithmetic then goes through smoothly and gives a wrong answer
with no warning. Fix the convention: \textbf{column $j$ is the input recipe
of sector $j$}, and column sums below $1$ mean each sector produces more
than it consumes.}

\wstep{Step 2. Write the balance equation.}
Each sector's output is used up as intermediate input plus final demand:
\[
  \mathbf x=A\mathbf x+\mathbf d .
\]
Collect the $\mathbf x$ terms --- carefully, since $\mathbf x$ is a vector:
\[
  \mathbf x-A\mathbf x=\mathbf d
  \;\Longrightarrow\;
  (I-A)\mathbf x=\mathbf d
  \;\Longrightarrow\;
  \mathbf x=(I-A)^{-1}\mathbf d .
\]
\wtrap{$\mathbf x-A\mathbf x=(I-A)\mathbf x$, not $(1-A)\mathbf x$. You
cannot subtract a matrix from the number $1$; the identity matrix is what
makes the factorisation legal, and forgetting it is why the inverse of the
wrong matrix gets computed.}

\wstep{Step 3. Form $I-A$ and invert it.}
\[
  I-A=\begin{pmatrix}0.8&-0.3\\-0.4&0.9\end{pmatrix},
  \qquad
  \det(I-A)=0.72-0.12=0.6\neq0 .
\]
\[
  (I-A)^{-1}=\frac{1}{0.6}\begin{pmatrix}0.9&0.3\\0.4&0.8\end{pmatrix}
  =\begin{pmatrix}1.5&0.5\\ \tfrac23&\tfrac43\end{pmatrix}.
\]
\wtrap{The $2\times2$ inverse swaps the \emph{diagonal} entries and negates
the \emph{off-diagonal} ones. Both off-diagonal entries of $I-A$ are already
negative, so negating them makes them positive --- which is why the Leontief
inverse comes out non-negative, as it must.}

\wstep{Step 4. Multiply out.}
\[
  \mathbf x=\frac{1}{0.6}\begin{pmatrix}0.9&0.3\\0.4&0.8\end{pmatrix}
             \begin{pmatrix}10\\20\end{pmatrix}
  =\frac{1}{0.6}\begin{pmatrix}9+6\\4+16\end{pmatrix}
  =\frac{1}{0.6}\begin{pmatrix}15\\20\end{pmatrix}
  =\boxed{\begin{pmatrix}25\\ 33.\overline{3}\end{pmatrix}}
\]

\wstep{Step 5. Check against the balance equation, not against the inverse.}
\[
  A\mathbf x=\begin{pmatrix}0.2\cdot25+0.3\cdot33.\overline{3}\\
                            0.4\cdot25+0.1\cdot33.\overline{3}\end{pmatrix}
  =\begin{pmatrix}5+10\\10+3.\overline{3}\end{pmatrix}
  =\begin{pmatrix}15\\13.\overline{3}\end{pmatrix},
\]
and $A\mathbf x+\mathbf d=(25,\,33.\overline{3})^{\!\top}=\mathbf x$.
\checkmark\ Note that gross output greatly exceeds final demand: producing
$10$ units for consumers requires $25$ units of gross output in sector~$1$,
because of the circular input requirements.

\wstep{Step 6. Productivity: is a non-negative solution guaranteed?}
The column sums of $A$ are $0.2+0.4=0.6$ and $0.3+0.1=0.4$, both $<1$: every
sector produces more than it uses up. That is the \textbf{Hawkins--Simon}
condition in its practical form, and it guarantees that $(I-A)^{-1}$ exists
and has only non-negative entries --- so a non-negative demand always yields
a non-negative, economically meaningful output.
\wtrap{A merely invertible $I-A$ is not enough. If a column sum exceeded
$1$, the algebra would still return a vector $\mathbf x$, but it could have
negative components --- a mathematically valid, economically impossible
answer. Check the column sums before trusting the output.}

\wstep{Step 7. Read the Leontief inverse as a table of multipliers.}
Entry $(i,j)$ of $(I-A)^{-1}$ is the gross output of sector~$i$ required per
unit of \emph{final} demand for good~$j$, direct and indirect effects
combined. Here $1.5$ units of good~$1$ are needed per unit of final demand
for good~$1$: one unit directly, half a unit through the supply chain. The
column sums, $1.5+\tfrac23\approx2.17$ and $0.5+\tfrac43\approx1.83$, are
the \textbf{output multipliers} --- exactly the numbers an input--output
study reports, and the reason a matrix inverse is worth computing here even
though solving one linear system would have been quicker.
\end{worked}

\section{Practice Problems}

\begin{problems}
  \item Leontief, $n = 2$. Let
        $A = \bigl(\begin{smallmatrix} 0.2 & 0.3 \\ 0.4 & 0.1 \end{smallmatrix}\bigr)$
        and $d = (10, 20)^{T}$. Compute the gross output vector $x$.
  \item For the same $A$, compute the Leontief inverse $(I - A)^{-1}$.
  \item Leontief, $n = 2$,
        $A = \bigl(\begin{smallmatrix} 0.5 & 0.25 \\ 0.4 & 0.5 \end{smallmatrix}\bigr)$,
        $d = (50, 60)^{T}$. Find $x$.
  \item Show that if every column of $A$ has sum $< 1$ (and $A \ge 0$
        entrywise) then $(I - A)$ is invertible. (Sketch: spectral radius
        argument or geometric series $\sum_{k \ge 0} A^{k}$ converges.)
  \item OLS with $n = 3, k = 2$. Data:
        $X = \begin{pmatrix} 1 & 1 \\ 1 & 2 \\ 1 & 3 \end{pmatrix}$,
        $y = (2, 2, 4)^{T}$. Compute $X^{T} X$, $X^{T} y$, and $\hat{\beta}$.
  \item For the same $X, y$, give the fitted values $\hat{y} = X \hat{\beta}$
        and the residual vector $y - \hat{y}$.
  \item OLS,
        $X = \begin{pmatrix} 1 & 0 \\ 1 & 1 \\ 1 & 2 \\ 1 & 3 \end{pmatrix}$,
        $y = (1, 1, 2, 2)^{T}$. Compute $\hat{\beta}$.
  \item Why does OLS require $X$ to have full column rank?
  \item Hessian classification. $f(x, y) = x^{2} + 4 x y + y^{2}$. Find and
        classify the critical points.
  \item $f(x, y) = x^{3} + y^{3} - 3 x y$. Find and classify all critical
        points.
  \item $f(x, y, z) = x^{2} + y^{2} + z^{2} - x y$. Find the critical point
        and classify.
  \item For $f(x, y) = \mathrm{e}^{x^{2} + y^{2}}$, classify the critical
        point at the origin.
  \item Markov.
        $P = \bigl(\begin{smallmatrix} 0.9 & 0.2 \\ 0.1 & 0.8 \end{smallmatrix}\bigr)$
        (column-stochastic). Find the stationary distribution.
  \item For
        $P = \bigl(\begin{smallmatrix} 0.6 & 0.3 \\ 0.4 & 0.7 \end{smallmatrix}\bigr)$,
        find the stationary distribution.
  \item For $P$ in problem 13, find $P^{n} v_{0}$ for $v_{0} = (1, 0)^{T}$
        as $n \to \infty$ (use the stationary distribution and the second
        eigenvalue).
  \item Portfolio. Two assets with covariance
        $\Sigma = \bigl(\begin{smallmatrix} 4 & 1 \\ 1 & 9 \end{smallmatrix}\bigr)$.
        Compute the variance of the equally weighted portfolio
        $w = (1/2, 1/2)^{T}$.
  \item For the same $\Sigma$, find the global minimum-variance portfolio
        (weights summing to $1$). \emph{Hint: minimise $w^{T} \Sigma w$
        subject to $w_{1} + w_{2} = 1$, e.g.\ set $w = (t, 1 - t)$ and
        minimise in $t$.}
  \item Three assets,
        $\Sigma = \begin{pmatrix} 2 & 0 & 0 \\ 0 & 3 & 0 \\ 0 & 0 & 5 \end{pmatrix}$.
        Compute the variance of the equally weighted portfolio
        $w = (1/3, 1/3, 1/3)^{T}$.
  \item For the same diagonal $\Sigma$, give a formula for the variance of
        an arbitrary portfolio $w = (w_{1}, w_{2}, w_{3})^{T}$.
  \item Show $X^{T} X$ is positive definite whenever $X$ has full column
        rank.
  \item OLS in a tiny one-regressor case.
        $X = (1, 2, 3, 4)^{T} \in \R^{4 \times 1}$ (a single column),
        $y = (2, 3, 5, 6)^{T}$. Compute
        $\hat{\beta} = (X^{T} X)^{-1} X^{T} y$ (a scalar).
  \item For an upper-triangular Markov matrix
        $P = \bigl(\begin{smallmatrix} 1 & 0.5 \\ 0 & 0.5 \end{smallmatrix}\bigr)$
        (column-stochastic), find the eigenvalues and stationary
        distribution.
  \item Verify the rule $\det(A^{-1}) = 1/\det A$ for $A = (I - A_{0})$ with
        $A_{0} = \bigl(\begin{smallmatrix} 0.2 & 0.3 \\ 0.4 & 0.1 \end{smallmatrix}\bigr)$
        from problem 1.
  \item For $f(x, y) = -x^{2} - y^{2} + 2 x + 4 y$, find the maximiser and
        the maximum value.
\end{problems}

\section{Solutions}

\begin{solutions}
  \item $I - A = \bigl(\begin{smallmatrix} 0.8 & -0.3 \\ -0.4 & 0.9 \end{smallmatrix}\bigr)$.
        $\det = 0.72 - 0.12 = 0.6$.
        $(I - A)^{-1} = \tfrac{1}{0.6}\bigl(\begin{smallmatrix} 0.9 & 0.3 \\ 0.4 & 0.8 \end{smallmatrix}\bigr) = \bigl(\begin{smallmatrix} 1.5 & 0.5 \\ 2/3 & 4/3 \end{smallmatrix}\bigr)$.
        \[
        x = (I - A)^{-1} d = \begin{pmatrix} 1.5 & 0.5 \\ 2/3 & 4/3 \end{pmatrix} \begin{pmatrix} 10 \\ 20 \end{pmatrix} = \begin{pmatrix} 15 + 10 \\ 20/3 + 80/3 \end{pmatrix} = \begin{pmatrix} 25 \\ 100/3 \end{pmatrix}.
        \]
        Check:
        $A x + d = \bigl(\begin{smallmatrix} 0.2 \cdot 25 + 0.3 \cdot 100/3 \\ 0.4 \cdot 25 + 0.1 \cdot 100/3 \end{smallmatrix}\bigr) + \bigl(\begin{smallmatrix} 10 \\ 20 \end{smallmatrix}\bigr) = \bigl(\begin{smallmatrix} 5 + 10 \\ 10 + 10/3 \end{smallmatrix}\bigr) + \bigl(\begin{smallmatrix} 10 \\ 20 \end{smallmatrix}\bigr) = \bigl(\begin{smallmatrix} 25 \\ 100/3 \end{smallmatrix}\bigr) = x$.
        Good.

  \item Computed in problem 1:
        $(I - A)^{-1} = \tfrac{1}{0.6}\bigl(\begin{smallmatrix} 0.9 & 0.3 \\ 0.4 & 0.8 \end{smallmatrix}\bigr) = \bigl(\begin{smallmatrix} 1.5 & 0.5 \\ 2/3 & 4/3 \end{smallmatrix}\bigr)$.

  \item $I - A = \bigl(\begin{smallmatrix} 0.5 & -0.25 \\ -0.4 & 0.5 \end{smallmatrix}\bigr)$.
        $\det = 0.25 - 0.1 = 0.15$.
        $(I - A)^{-1} = \tfrac{1}{0.15} \bigl(\begin{smallmatrix} 0.5 & 0.25 \\ 0.4 & 0.5 \end{smallmatrix}\bigr) = \tfrac{1}{0.15}\bigl(\begin{smallmatrix} 0.5 & 0.25 \\ 0.4 & 0.5 \end{smallmatrix}\bigr)$.
        \[
        x = (I - A)^{-1} d = \tfrac{1}{0.15}\begin{pmatrix} 0.5 \cdot 50 + 0.25 \cdot 60 \\ 0.4 \cdot 50 + 0.5 \cdot 60 \end{pmatrix} = \tfrac{1}{0.15}\begin{pmatrix} 40 \\ 50 \end{pmatrix} = \begin{pmatrix} 800/3 \\ 1000/3 \end{pmatrix}.
        \]
        Check: $A x + d$, first entry:
        $0.5 \cdot 800/3 + 0.25 \cdot 1000/3 + 50 = 400/3 + 250/3 + 150/3 = 800/3$.
        Good.

  \item If column sums of $A$ are all $< 1$ (and $A \ge 0$), the matrix $A$
        has spectral radius $\rho(A) < 1$ (by Perron-Frobenius, since the
        spectral radius of a nonnegative matrix is bounded by the maximal
        column sum). Hence the geometric series
        $\sum_{k = 0}^{\infty} A^{k}$ converges and equals $(I - A)^{-1}$;
        in particular $I - A$ is invertible.

  \item $X^{T} X = \begin{pmatrix} 1 & 1 & 1 \\ 1 & 2 & 3 \end{pmatrix} \begin{pmatrix} 1 & 1 \\ 1 & 2 \\ 1 & 3 \end{pmatrix} = \bigl(\begin{smallmatrix} 3 & 6 \\ 6 & 14 \end{smallmatrix}\bigr)$.
        With $X^{T} y = (\sum y_{i}, \sum x_{i} y_{i})^{T}$ where $x_{i}$ is
        column 2 of $X$: $\sum y_{i} = 2 + 2 + 4 = 8$ and
        $\sum x_{i} y_{i} = 1 \cdot 2 + 2 \cdot 2 + 3 \cdot 4 = 2 + 4 + 12 = 18$.
        So $X^{T} y = (8, 18)^{T}$.
        $\det(X^{T} X) = 42 - 36 = 6$.
        $(X^{T} X)^{-1} = \tfrac{1}{6}\bigl(\begin{smallmatrix} 14 & -6 \\ -6 & 3 \end{smallmatrix}\bigr)$.
        \[
        \hat{\beta} = \tfrac{1}{6}\begin{pmatrix} 14 \cdot 8 - 6 \cdot 18 \\ -6 \cdot 8 + 3 \cdot 18 \end{pmatrix} = \tfrac{1}{6}\begin{pmatrix} 112 - 108 \\ -48 + 54 \end{pmatrix} = \tfrac{1}{6}\begin{pmatrix} 4 \\ 6 \end{pmatrix} = \begin{pmatrix} 2/3 \\ 1 \end{pmatrix}.
        \]

  \item $\hat{y} = X \hat{\beta} = \begin{pmatrix} 1 \\ 1 \\ 1 \end{pmatrix}(2/3) + \begin{pmatrix} 1 \\ 2 \\ 3 \end{pmatrix}(1) = \begin{pmatrix} 2/3 + 1 \\ 2/3 + 2 \\ 2/3 + 3 \end{pmatrix} = \begin{pmatrix} 5/3 \\ 8/3 \\ 11/3 \end{pmatrix}$.
        Residuals
        $y - \hat{y} = (2 - 5/3, 2 - 8/3, 4 - 11/3)^{T} = (1/3, -2/3, 1/3)^{T}$.
        (Note they sum to $0$ and are orthogonal to each column of $X$, as
        required.)

  \item $X^{T} X = \bigl(\begin{smallmatrix} 4 & 6 \\ 6 & 14 \end{smallmatrix}\bigr)$,
        where $\sum x_{i}^{2} = 0 + 1 + 4 + 9 = 14$, $\sum x_{i} = 6$.
        $X^{T} y = (\sum y_{i}, \sum x_{i} y_{i})^{T} = (6, 0 + 1 + 4 + 6)^{T} = (6, 11)^{T}$.
        $\det(X^{T} X) = 56 - 36 = 20$.
        $(X^{T} X)^{-1} = \tfrac{1}{20}\bigl(\begin{smallmatrix} 14 & -6 \\ -6 & 4 \end{smallmatrix}\bigr)$.
        So
        \[
        \hat{\beta} = \tfrac{1}{20}\begin{pmatrix} 14 \cdot 6 - 6 \cdot 11 \\ -6 \cdot 6 + 4 \cdot 11 \end{pmatrix} = \tfrac{1}{20}\begin{pmatrix} 84 - 66 \\ -36 + 44 \end{pmatrix} = \tfrac{1}{20}\begin{pmatrix} 18 \\ 8 \end{pmatrix} = \begin{pmatrix} 0.9 \\ 0.4 \end{pmatrix}.
        \]

  \item If $X$ lacks full column rank, two columns (regressors) are linearly
        dependent. Then $X^{T} X$ is singular ($\det = 0$,
        $\mathrm{rk}(X^{T} X) = \mathrm{rk}(X) < k$), so the inverse
        $(X^{T} X)^{-1}$ does not exist and $\hat{\beta}$ is not uniquely
        defined. Geometrically: the projection of $y$ onto
        $\mathrm{col}(X)$ is still well-defined, but multiple coefficient
        vectors $\beta$ produce the same projection.

  \item $f_{x} = 2 x + 4 y, f_{y} = 4 x + 2 y$.
        $f_{x} = f_{y} = 0 \Rightarrow x = y = 0$.
        $H = \bigl(\begin{smallmatrix} 2 & 4 \\ 4 & 2 \end{smallmatrix}\bigr)$,
        $D_{2} = 4 - 16 = -12 < 0$: indefinite. Origin is a saddle.

  \item $f_{x} = 3 x^{2} - 3 y = 0 \Rightarrow y = x^{2}$.
        $f_{y} = 3 y^{2} - 3 x = 0 \Rightarrow x = y^{2}$. Combine:
        $x = x^{4}$, so $x(x^{3} - 1) = 0$, giving $x = 0$ or $x = 1$.
        Critical points $(0, 0)$ and $(1, 1)$.
        $H = \bigl(\begin{smallmatrix} 6 x & -3 \\ -3 & 6 y \end{smallmatrix}\bigr)$.
        At $(0, 0)$:
        $H = \bigl(\begin{smallmatrix} 0 & -3 \\ -3 & 0 \end{smallmatrix}\bigr)$,
        $D_{2} = -9 < 0$: indefinite, saddle. At $(1, 1)$:
        $H = \bigl(\begin{smallmatrix} 6 & -3 \\ -3 & 6 \end{smallmatrix}\bigr)$,
        $D_{1} = 6, D_{2} = 27 > 0$: pd, strict local min.

  \item $f_{x} = 2 x - y, f_{y} = 2 y - x, f_{z} = 2 z$. All zero iff
        $x = y = z = 0$.
        $H = \begin{pmatrix} 2 & -1 & 0 \\ -1 & 2 & 0 \\ 0 & 0 & 2 \end{pmatrix}$,
        $D_{1} = 2, D_{2} = 3, D_{3} = 2 \cdot 3 = 6$: pd. Origin is a
        strict local (in fact global) min.

  \item $f_{x} = 2 x \mathrm{e}^{x^{2} + y^{2}}, f_{y} = 2 y \mathrm{e}^{x^{2} + y^{2}}$.
        Both zero iff $x = y = 0$. At origin,
        $H = \bigl(\begin{smallmatrix} 2 & 0 \\ 0 & 2 \end{smallmatrix}\bigr) \cdot \mathrm{e}^{0} = 2 I$
        (after differentiating again at $0$; cross terms vanish). pd:
        strict local min.

  \item $(P - I)\pi = 0$:
        $\bigl(\begin{smallmatrix} -0.1 & 0.2 \\ 0.1 & -0.2 \end{smallmatrix}\bigr)\pi = 0$.
        So $-0.1 \pi_{1} + 0.2 \pi_{2} = 0 \Rightarrow \pi_{1} = 2 \pi_{2}$.
        With $\pi_{1} + \pi_{2} = 1$: $\pi = (2/3, 1/3)^{T}$.

  \item $(P - I)\pi = 0$:
        $\bigl(\begin{smallmatrix} -0.4 & 0.3 \\ 0.4 & -0.3 \end{smallmatrix}\bigr)\pi = 0 \Rightarrow 4 \pi_{1} = 3 \pi_{2}$,
        i.e.\ $\pi_{1} = 3/7, \pi_{2} = 4/7$.

  \item Eigenvalues of $P$ from problem 13:
        $\lambda^{2} - 1.7 \lambda + (0.72 - 0.02) = \lambda^{2} - 1.7 \lambda + 0.7 = (\lambda - 1)(\lambda - 0.7)$.
        The second eigenvalue $0.7 < 1$, so $P^{n}$ converges.
        $\lim_{n \to \infty} P^{n} v_{0}$ is the stationary distribution
        $(2/3, 1/3)^{T}$.

  \item $w^{T} \Sigma w = (1/2, 1/2)\bigl(\begin{smallmatrix} 4 & 1 \\ 1 & 9 \end{smallmatrix}\bigr)\bigl(\begin{smallmatrix} 1/2 \\ 1/2 \end{smallmatrix}\bigr) = (1/2, 1/2)\bigl(\begin{smallmatrix} 5/2 \\ 5 \end{smallmatrix}\bigr) = 5/4 + 5/2 = 15/4$.

  \item Set $w = (t, 1 - t)$:
        $w^{T} \Sigma w = 4 t^{2} + 2 \cdot 1 \cdot t(1 - t) + 9(1 - t)^{2} = 4 t^{2} + 2 t - 2 t^{2} + 9 - 18 t + 9 t^{2} = 11 t^{2} - 16 t + 9$.
        Differentiate: $22 t - 16 = 0 \Rightarrow t = 8/11$. So
        $w^{*} = (8/11, 3/11)^{T}$. Minimum variance
        $= 11 \cdot 64/121 - 16 \cdot 8/11 + 9 = 64/11 - 128/11 + 99/11 = 35/11$.

  \item $w^{T} \Sigma w = (1/9)(2) + (1/9)(3) + (1/9)(5) = 10/9$.

  \item Diagonal $\Sigma$:
        $w^{T} \Sigma w = 2 w_{1}^{2} + 3 w_{2}^{2} + 5 w_{3}^{2}$.

  \item For $x \neq 0$, $x^{T}(X^{T} X)x = \|X x\|^{2} \ge 0$. If
        $\|X x\|^{2} = 0$ then $X x = 0$. Since $X$ has full column rank,
        the only solution is $x = 0$. So $x^{T}(X^{T} X)x > 0$ for
        $x \neq 0$: pd.

  \item $X^{T} X = 1 + 4 + 9 + 16 = 30$.
        $X^{T} y = 1 \cdot 2 + 2 \cdot 3 + 3 \cdot 5 + 4 \cdot 6 = 2 + 6 + 15 + 24 = 47$.
        $\hat{\beta} = 47/30$.

  \item Upper-triangular: eigenvalues $\lambda = 1, 0.5$ (the diagonal).
        For $\lambda = 1$:
        $(P - I) = \bigl(\begin{smallmatrix} 0 & 0.5 \\ 0 & -0.5 \end{smallmatrix}\bigr)$,
        so $\pi_{2} = 0$, $\pi_{1}$ free. Normalised: $\pi = (1, 0)^{T}$.
        (State $1$ is absorbing.)

  \item From problem 1, $\det(I - A_{0}) = 0.6$, so
        $\det((I - A_{0})^{-1}) = 1/0.6 = 5/3$. Directly: from problem 1,
        $(I - A_{0})^{-1} = \bigl(\begin{smallmatrix} 1.5 & 0.5 \\ 2/3 & 4/3 \end{smallmatrix}\bigr)$,
        $\det = 1.5 \cdot 4/3 - 0.5 \cdot 2/3 = 2 - 1/3 = 5/3$. Match.

  \item $f_{x} = -2 x + 2 = 0 \Rightarrow x = 1$.
        $f_{y} = -2 y + 4 = 0 \Rightarrow y = 2$. $H = -2 I$, nd: strict
        (and global) max. $f(1, 2) = -1 - 4 + 2 + 8 = 5$.
\end{solutions}

\section{Additional Problems}

\begin{problems}
  \item Leontief, $n = 2$. Let
        $A = \bigl(\begin{smallmatrix} 0.1 & 0.2 \\ 0.3 & 0.4 \end{smallmatrix}\bigr)$
        and $d = (8, 12)^{T}$. Compute the gross output vector $x$.
  \item For the same $A$, give the Leontief inverse $(I - A)^{-1}$
        explicitly.
  \item Leontief, $n = 2$.
        $A = \bigl(\begin{smallmatrix} 0.2 & 0.1 \\ 0.5 & 0.3 \end{smallmatrix}\bigr)$,
        $d = (4, 7)^{T}$. Find $x$.
  \item For a Leontief economy with
        $A = \bigl(\begin{smallmatrix} 0.5 & 0.2 \\ 0.4 & 0.5 \end{smallmatrix}\bigr)$,
        check that the spectral radius of $A$ is $< 1$ (so productivity
        holds) and compute $(I - A)^{-1}$.
  \item OLS with $n = 4, k = 2$. Data:
        $X = \begin{pmatrix} 1 & 1 \\ 1 & 2 \\ 1 & 3 \\ 1 & 4 \end{pmatrix}$,
        $y = (3, 5, 7, 10)^{T}$. Compute $X^{T} X$, $X^{T} y$, and
        $\hat{\beta}$.
  \item For the same $X, y$, give the fitted values $\hat{y}$ and residuals
        $y - \hat{y}$.
  \item OLS,
        $X = \begin{pmatrix} 1 & -1 \\ 1 & 0 \\ 1 & 1 \\ 1 & 2 \end{pmatrix}$,
        $y = (0, 1, 3, 4)^{T}$. Compute $\hat{\beta}$.
  \item Hessian classification.
        $f(x, y) = 2 x^{2} + y^{2} + 4 x y - 6 x$. Find and classify
        critical points.
  \item $f(x, y) = x^{2} - y^{2} + 4 x + 6 y$. Find and classify the
        critical point.
  \item $f(x, y, z) = x^{2} + 2 y^{2} + 3 z^{2} - 2 x y - 4 y z + 5$. Find
        the critical point and classify.
  \item For $f(x, y) = \ln(1 + x^{2} + y^{2})$, classify the critical point
        at the origin.
  \item Markov.
        $P = \bigl(\begin{smallmatrix} 0.7 & 0.4 \\ 0.3 & 0.6 \end{smallmatrix}\bigr)$
        (column-stochastic). Find the stationary distribution.
  \item For
        $P = \bigl(\begin{smallmatrix} 0.5 & 0.2 \\ 0.5 & 0.8 \end{smallmatrix}\bigr)$,
        find the stationary distribution.
  \item For $P$ in problem 12, compute the eigenvalues and use them to argue
        $\lim_{n \to \infty} P^{n} v_{0}$ for any probability vector $v_{0}$
        equals the stationary distribution.
  \item Portfolio. Two assets with covariance
        $\Sigma = \bigl(\begin{smallmatrix} 9 & 2 \\ 2 & 4 \end{smallmatrix}\bigr)$.
        Compute the variance of $w = (1/3, 2/3)^{T}$.
  \item For the same $\Sigma$, find the global minimum-variance portfolio.
  \item Three assets,
        $\Sigma = \begin{pmatrix} 4 & 1 & 0 \\ 1 & 5 & 0 \\ 0 & 0 & 9 \end{pmatrix}$.
        Compute the variance of the equally weighted portfolio
        $w = (1/3, 1/3, 1/3)^{T}$.
  \item For the same $\Sigma$, classify $\Sigma$ (pd, psd, indefinite?). Use
        principal minors.
  \item OLS, single regressor through origin.
        $X = (1, 2, 3)^{T} \in \R^{3 \times 1}$, $y = (1, 3, 4)^{T}$.
        Compute $\hat{\beta}$ (a scalar).
  \item Markov.
        $P = \begin{pmatrix} 0.5 & 0.1 & 0.2 \\ 0.3 & 0.6 & 0.3 \\ 0.2 & 0.3 & 0.5 \end{pmatrix}$
        (column-stochastic). Verify $\pi = (1/3, 1/3, 1/3)^{T}$ \emph{or}
        find the stationary distribution.
  \item For $f(x, y) = x^{2} + x y + y^{2} + 2 x - 4 y$, find the critical
        point and the optimum value.
  \item Verify for the OLS problem in Additional Problem 5 that the residual
        vector is orthogonal to each column of $X$.
  \item Leontief with productivity check: for
        $A = \bigl(\begin{smallmatrix} 0.3 & 0.6 \\ 0.4 & 0.5 \end{smallmatrix}\bigr)$,
        determine whether the productivity condition (column sums $<1$)
        holds. If yes, find $x$ for $d = (10, 10)^{T}$; if no, explain.
  \item Portfolio with constraint. Two assets,
        $\Sigma = \bigl(\begin{smallmatrix} 1 & 0.5 \\ 0.5 & 4 \end{smallmatrix}\bigr)$.
        Find the global minimum-variance portfolio and its variance.
\end{problems}

\section{Additional Solutions}

\begin{solutions}
  \item $I - A = \bigl(\begin{smallmatrix} 0.9 & -0.2 \\ -0.3 & 0.6 \end{smallmatrix}\bigr)$.
        $\det = 0.54 - 0.06 = 0.48$.
        $(I - A)^{-1} = \tfrac{1}{0.48}\bigl(\begin{smallmatrix} 0.6 & 0.2 \\ 0.3 & 0.9 \end{smallmatrix}\bigr) = \bigl(\begin{smallmatrix} 5/4 & 5/12 \\ 5/8 & 15/8 \end{smallmatrix}\bigr)$.
        \[
        x = (I - A)^{-1} d = \tfrac{1}{0.48}\begin{pmatrix} 0.6 & 0.2 \\ 0.3 & 0.9 \end{pmatrix}\begin{pmatrix} 8 \\ 12 \end{pmatrix} = \tfrac{1}{0.48}\begin{pmatrix} 4.8 + 2.4 \\ 2.4 + 10.8 \end{pmatrix} = \tfrac{1}{0.48}\begin{pmatrix} 7.2 \\ 13.2 \end{pmatrix} = \begin{pmatrix} 15 \\ 27.5 \end{pmatrix}.
        \]
        Verify:
        $A x + d = (0.1 \cdot 15 + 0.2 \cdot 27.5, 0.3 \cdot 15 + 0.4 \cdot 27.5)^{T} + (8, 12)^{T} = (1.5 + 5.5, 4.5 + 11)^{T} + (8, 12)^{T} = (7, 15.5)^{T} + (8, 12)^{T} = (15, 27.5)^{T}$.
        Good.

  \item From problem 1:
        $(I - A)^{-1} = \bigl(\begin{smallmatrix} 5/4 & 5/12 \\ 5/8 & 15/8 \end{smallmatrix}\bigr)$.

  \item $I - A = \bigl(\begin{smallmatrix} 0.8 & -0.1 \\ -0.5 & 0.7 \end{smallmatrix}\bigr)$.
        $\det = 0.56 - 0.05 = 0.51$.
        $(I - A)^{-1} = \tfrac{1}{0.51}\bigl(\begin{smallmatrix} 0.7 & 0.1 \\ 0.5 & 0.8 \end{smallmatrix}\bigr)$.
        \[
        x = \tfrac{1}{0.51}\begin{pmatrix} 0.7 \cdot 4 + 0.1 \cdot 7 \\ 0.5 \cdot 4 + 0.8 \cdot 7 \end{pmatrix} = \tfrac{1}{0.51}\begin{pmatrix} 3.5 \\ 7.6 \end{pmatrix} = \begin{pmatrix} 350/51 \\ 760/51 \end{pmatrix} \approx \begin{pmatrix} 6.86 \\ 14.90 \end{pmatrix}.
        \]

  \item Column sums: $0.5 + 0.4 = 0.9 < 1$ and $0.2 + 0.5 = 0.7 < 1$, so
        spectral radius bound $\rho(A) \le \max$ col sum $= 0.9 < 1$:
        productive.
        $I - A = \bigl(\begin{smallmatrix} 0.5 & -0.2 \\ -0.4 & 0.5 \end{smallmatrix}\bigr)$.
        $\det = 0.25 - 0.08 = 0.17$.
        $(I - A)^{-1} = \tfrac{1}{0.17}\bigl(\begin{smallmatrix} 0.5 & 0.2 \\ 0.4 & 0.5 \end{smallmatrix}\bigr)$.

  \item $X^{T} X = \bigl(\begin{smallmatrix} n & \sum x_{i} \\ \sum x_{i} & \sum x_{i}^{2} \end{smallmatrix}\bigr) = \bigl(\begin{smallmatrix} 4 & 10 \\ 10 & 30 \end{smallmatrix}\bigr)$,
        since $\sum x_{i} = 10, \sum x_{i}^{2} = 1 + 4 + 9 + 16 = 30$.
        $X^{T} y = (\sum y_{i}, \sum x_{i} y_{i})^{T} = (3 + 5 + 7 + 10, 3 + 10 + 21 + 40)^{T} = (25, 74)^{T}$.
        $\det(X^{T} X) = 120 - 100 = 20$.
        $(X^{T} X)^{-1} = \tfrac{1}{20}\bigl(\begin{smallmatrix} 30 & -10 \\ -10 & 4 \end{smallmatrix}\bigr)$.
        So
        \[
        \hat{\beta} = \tfrac{1}{20}\begin{pmatrix} 30 \cdot 25 - 10 \cdot 74 \\ -10 \cdot 25 + 4 \cdot 74 \end{pmatrix} = \tfrac{1}{20}\begin{pmatrix} 750 - 740 \\ -250 + 296 \end{pmatrix} = \tfrac{1}{20}\begin{pmatrix} 10 \\ 46 \end{pmatrix} = \begin{pmatrix} 0.5 \\ 2.3 \end{pmatrix}.
        \]

  \item $\hat{y}_{i} = 0.5 + 2.3 x_{i}$:
        $\hat{y} = (0.5 + 2.3, 0.5 + 4.6, 0.5 + 6.9, 0.5 + 9.2)^{T} = (2.8, 5.1, 7.4, 9.7)^{T}$.
        Residuals
        $y - \hat{y} = (3 - 2.8, 5 - 5.1, 7 - 7.4, 10 - 9.7)^{T} = (0.2, -0.1, -0.4, 0.3)^{T}$.
        Sum $= 0$: good.

  \item $\sum x_{i} = -1 + 0 + 1 + 2 = 2$,
        $\sum x_{i}^{2} = 1 + 0 + 1 + 4 = 6$,
        $\sum y_{i} = 0 + 1 + 3 + 4 = 8$,
        $\sum x_{i} y_{i} = 0 + 0 + 3 + 8 = 11$.
        $X^{T} X = \bigl(\begin{smallmatrix} 4 & 2 \\ 2 & 6 \end{smallmatrix}\bigr)$,
        $\det = 20$.
        $(X^{T} X)^{-1} = \tfrac{1}{20}\bigl(\begin{smallmatrix} 6 & -2 \\ -2 & 4 \end{smallmatrix}\bigr)$.
        So
        \[
        \hat{\beta} = \tfrac{1}{20}\begin{pmatrix} 6 \cdot 8 - 2 \cdot 11 \\ -2 \cdot 8 + 4 \cdot 11 \end{pmatrix} = \tfrac{1}{20}\begin{pmatrix} 26 \\ 28 \end{pmatrix} = \begin{pmatrix} 1.3 \\ 1.4 \end{pmatrix}.
        \]

  \item $f_{x} = 4 x + 4 y - 6 = 0$,
        $f_{y} = 2 y + 4 x = 0 \Rightarrow y = -2 x$. Substitute:
        $4 x - 8 x - 6 = 0 \Rightarrow x = -3/2, y = 3$.
        $H = \bigl(\begin{smallmatrix} 4 & 4 \\ 4 & 2 \end{smallmatrix}\bigr)$,
        $D_{2} = 8 - 16 = -8 < 0$: indefinite. Saddle at $(-3/2, 3)$.

  \item $f_{x} = 2 x + 4 = 0 \Rightarrow x = -2$.
        $f_{y} = -2 y + 6 = 0 \Rightarrow y = 3$.
        $H = \bigl(\begin{smallmatrix} 2 & 0 \\ 0 & -2 \end{smallmatrix}\bigr)$,
        indefinite (eigenvalues $2, -2$). Saddle at $(-2, 3)$.

  \item $f_{x} = 2 x - 2 y = 0 \Rightarrow x = y$.
        $f_{y} = 4 y - 2 x - 4 z = 0 \Rightarrow 4 y - 2 y - 4 z = 0 \Rightarrow z = y/2$.
        $f_{z} = 6 z - 4 y = 0 \Rightarrow z = 2 y/3$. So
        $y/2 = 2 y/3 \Rightarrow 3 y = 4 y \Rightarrow y = 0$. So
        $x = y = z = 0$.
        $H = \begin{pmatrix} 2 & -2 & 0 \\ -2 & 4 & -4 \\ 0 & -4 & 6 \end{pmatrix}$.
        $D_{1} = 2 > 0$, $D_{2} = 8 - 4 = 4 > 0$,
        $D_{3} = 2(24 - 16) - (-2)(-12 - 0) + 0 = 16 - 24 = -8$. Signs
        $+, +, -$: not pd, not nd, so indefinite. Origin is a saddle.

  \item $f_{x} = \tfrac{2 x}{1 + x^{2} + y^{2}}$,
        $f_{y} = \tfrac{2 y}{1 + x^{2} + y^{2}}$. Both zero at origin. At
        origin,
        $H = \bigl(\begin{smallmatrix} 2 & 0 \\ 0 & 2 \end{smallmatrix}\bigr)$
        (second derivatives evaluated at $(0,0)$ give
        $f_{xx} = 2, f_{yy} = 2, f_{xy} = 0$). pd: strict local min.

  \item $(P - I)\pi = 0$:
        $\bigl(\begin{smallmatrix} -0.3 & 0.4 \\ 0.3 & -0.4 \end{smallmatrix}\bigr)\pi = 0 \Rightarrow 0.3 \pi_{1} = 0.4 \pi_{2}$,
        so $\pi_{1} = \tfrac{4}{3} \pi_{2}$. With $\pi_{1} + \pi_{2} = 1$:
        $\pi_{2} = 3/7, \pi_{1} = 4/7$.

  \item $(P - I)\pi = 0$:
        $\bigl(\begin{smallmatrix} -0.5 & 0.2 \\ 0.5 & -0.2 \end{smallmatrix}\bigr)\pi = 0 \Rightarrow 5 \pi_{1} = 2 \pi_{2}$,
        so $\pi_{1} = \tfrac{2}{7}, \pi_{2} = \tfrac{5}{7}$. Verify:
        $P \pi = (0.5 \cdot 2/7 + 0.2 \cdot 5/7, 0.5 \cdot 2/7 + 0.8 \cdot 5/7)^{T} = (1/7 + 1/7, 1/7 + 4/7)^{T} = (2/7, 5/7)^{T} = \pi$.
        Good.

  \item $p(\lambda) = (0.7 - \lambda)(0.6 - \lambda) - 0.12 = \lambda^{2} - 1.3 \lambda + 0.42 - 0.12 = \lambda^{2} - 1.3 \lambda + 0.30 = (\lambda - 1)(\lambda - 0.3)$.
        Since the second eigenvalue $0.3 < 1$, $P^{n} v_{0} \to $
        stationary distribution as $n \to \infty$, which is
        $(4/7, 3/7)^{T}$ from problem 12.

  \item $w^{T} \Sigma w = (1/3, 2/3)\bigl(\begin{smallmatrix} 9 & 2 \\ 2 & 4 \end{smallmatrix}\bigr)\bigl(\begin{smallmatrix} 1/3 \\ 2/3 \end{smallmatrix}\bigr) = (1/3, 2/3)\bigl(\begin{smallmatrix} 3 + 4/3 \\ 2/3 + 8/3 \end{smallmatrix}\bigr) = (1/3, 2/3)\bigl(\begin{smallmatrix} 13/3 \\ 10/3 \end{smallmatrix}\bigr) = 13/9 + 20/9 = 33/9 = 11/3$.

  \item Set $w = (t, 1 - t)$.
        $w^{T} \Sigma w = 9 t^{2} + 4 t(1 - t) + 4(1 - t)^{2} = 9 t^{2} + 4 t - 4 t^{2} + 4 - 8 t + 4 t^{2} = 9 t^{2} - 4 t + 4$.
        Differentiate: $18 t - 4 = 0 \Rightarrow t = 2/9$. So
        $w^{*} = (2/9, 7/9)^{T}$. Variance
        $= 9 \cdot 4/81 - 4 \cdot 2/9 + 4 = 4/9 - 8/9 + 36/9 = 32/9$.

  \item Compute
        $\Sigma w = \tfrac{1}{3}(4 + 1, 1 + 5, 9)^{T} = (5/3, 6/3, 9/3)^{T} = (5/3, 2, 3)^{T}$.
        Then
        $w^{T} \Sigma w = (1/3)(5/3 + 2 + 3) = (1/3)(5/3 + 5) = (1/3)(20/3) = 20/9$.

  \item Block diagonal. Upper $2 \times 2$: $D_{1} = 4 > 0$,
        $D_{2} = 20 - 1 = 19 > 0$: pd. Lower $1 \times 1$ block: $9 > 0$.
        So $\Sigma$ is pd overall.

  \item $X^{T} X = 1 + 4 + 9 = 14$. $X^{T} y = 1 + 6 + 12 = 19$.
        $\hat{\beta} = 19/14$.

  \item Test $\pi = (1/3, 1/3, 1/3)^{T}$.
        $P \pi = \tfrac{1}{3}(0.5 + 0.1 + 0.2, 0.3 + 0.6 + 0.3, 0.2 + 0.3 + 0.5)^{T} = \tfrac{1}{3}(0.8, 1.2, 1.0)^{T} = (0.267, 0.4, 0.333)^{T} \neq \pi$.
        So uniform is not stationary. Solve $(P - I)\pi = 0$ with
        $\pi_{1} + \pi_{2} + \pi_{3} = 1$:
        \[
        P - I = \begin{pmatrix} -0.5 & 0.1 & 0.2 \\ 0.3 & -0.4 & 0.3 \\ 0.2 & 0.3 & -0.5 \end{pmatrix}.
        \]
        Take row 1 and 2:
        $-0.5 \pi_{1} + 0.1 \pi_{2} + 0.2 \pi_{3} = 0$,
        $0.3 \pi_{1} - 0.4 \pi_{2} + 0.3 \pi_{3} = 0$. Multiply by $10$:
        $-5 \pi_{1} + \pi_{2} + 2 \pi_{3} = 0$,
        $3 \pi_{1} - 4 \pi_{2} + 3 \pi_{3} = 0$. Add $4 \times$ first to
        second:
        $-20 \pi_{1} + 4 \pi_{2} + 8 \pi_{3} + 3 \pi_{1} - 4 \pi_{2} + 3 \pi_{3} = 0 \Rightarrow -17 \pi_{1} + 11 \pi_{3} = 0 \Rightarrow \pi_{3} = \tfrac{17}{11} \pi_{1}$.
        From first:
        $\pi_{2} = 5 \pi_{1} - 2 \pi_{3} = 5 \pi_{1} - \tfrac{34}{11} \pi_{1} = \tfrac{55 - 34}{11} \pi_{1} = \tfrac{21}{11} \pi_{1}$.
        So
        $\pi = \pi_{1}(1, 21/11, 17/11)^{T} = \tfrac{\pi_{1}}{11}(11, 21, 17)^{T}$.
        Sum: $11 + 21 + 17 = 49$, set
        $\pi_{1}/11 \cdot 49 = 1 \Rightarrow \pi = \tfrac{1}{49}(11, 21, 17)^{T}$.
        Verify column sums of $P$: $0.5 + 0.3 + 0.2 = 1$,
        $0.1 + 0.6 + 0.3 = 1$, $0.2 + 0.3 + 0.5 = 1$. Good. Check $P \pi$:
        first entry
        $= \tfrac{1}{49}(0.5 \cdot 11 + 0.1 \cdot 21 + 0.2 \cdot 17) = \tfrac{1}{49}(5.5 + 2.1 + 3.4) = \tfrac{11}{49}$.
        Match.

  \item $f_{x} = 2 x + y + 2 = 0$, $f_{y} = x + 2 y - 4 = 0$. Solve: from
        first, $y = -2 x - 2$; sub:
        $x + 2(-2 x - 2) - 4 = 0 \Rightarrow -3 x - 8 = 0 \Rightarrow x = -8/3, y = -2(-8/3) - 2 = 16/3 - 6/3 = 10/3$.
        $H = \bigl(\begin{smallmatrix} 2 & 1 \\ 1 & 2 \end{smallmatrix}\bigr)$,
        pd: local (and global) min. Value:
        $f(-8/3, 10/3) = 64/9 + (-8/3)(10/3) + 100/9 + 2(-8/3) - 4(10/3) = 64/9 - 80/9 + 100/9 - 16/3 - 40/3 = 84/9 - 56/3 = 84/9 - 168/9 = -84/9 = -28/3$.

  \item Residuals $r = (0.2, -0.1, -0.4, 0.3)^{T}$. Col 1 of $X$ is
        $\mathbf{1}$: $\mathbf{1}^{T} r = 0.2 - 0.1 - 0.4 + 0.3 = 0$. Col 2
        of $X$ is $(1, 2, 3, 4)^{T}$:
        $1 \cdot 0.2 + 2 \cdot (-0.1) + 3 \cdot (-0.4) + 4 \cdot 0.3 = 0.2 - 0.2 - 1.2 + 1.2 = 0$.
        Good.

  \item Column sums: $0.3 + 0.4 = 0.7 < 1$ and
        $0.6 + 0.5 = 1.1 \not< 1$. Productivity fails in the column-sum
        sense. However we can still check directly:
        $\det(I - A) = (0.7)(0.5) - (-0.6)(-0.4) = 0.35 - 0.24 = 0.11 \neq 0$,
        so $I - A$ is invertible.
        $(I - A)^{-1} = \tfrac{1}{0.11}\bigl(\begin{smallmatrix} 0.5 & 0.6 \\ 0.4 & 0.7 \end{smallmatrix}\bigr)$.
        $x = \tfrac{1}{0.11}(0.5 \cdot 10 + 0.6 \cdot 10, 0.4 \cdot 10 + 0.7 \cdot 10)^{T} = \tfrac{1}{0.11}(11, 11)^{T} = (100, 100)^{T}$.
        Productivity in spectral-radius sense does hold here (since the
        inverse exists with positive entries); the column-sum bound is
        sufficient but not necessary.

  \item Set $w = (t, 1 - t)$.
        $w^{T} \Sigma w = t^{2} + 2 \cdot 0.5 \cdot t(1 - t) + 4(1 - t)^{2} = t^{2} + t - t^{2} + 4 - 8 t + 4 t^{2} = 4 t^{2} - 7 t + 4$.
        Differentiate: $8 t - 7 = 0 \Rightarrow t = 7/8$. So
        $w^{*} = (7/8, 1/8)^{T}$. Variance
        $= 4(49/64) - 7(7/8) + 4 = 49/16 - 49/8 + 64/16 = 49/16 - 98/16 + 64/16 = 15/16$.
\end{solutions}
